\documentclass[a4paper,12pt]{book}

%Pacchetti utili
%************
%* packages *
%************

%Codifica
\usepackage[utf8]{inputenc}

%Lingua, sostituire italian con english nel caso in cui la tesi sia scritta in inglese
\usepackage[italian]{babel}

%Pacchetto per definire layout di pagina
\usepackage{fancyhdr}
\usepackage{sectsty}
\usepackage[left=3cm, right=3cm, bottom=3cm]{geometry}

%Spazia linee all'interno del documento
\usepackage{setspace}

%Listati di codice
\usepackage{verbatim}
\usepackage{listings}

%Didascalie immagini
\usepackage[hang,small,sf,font=scriptsize, labelfont=bf]{caption}
\usepackage{subcaption}

%Inclusione immagini
\usepackage{graphicx}

%Impostazioni note a piè pagina
\usepackage[stable]{footmisc}

%Citazioni e riferimenti a label
\usepackage{cite}
\usepackage[english]{varioref}

%Colori
\usepackage[usenames]{color}
\usepackage{xcolor}
\usepackage{colortbl}

%Crea link ipertestuali
\usepackage[hidelinks]{hyperref}

%Formattazione url
\usepackage{url}

%Inserimento formule
\usepackage{amsmath}
\usepackage{mathrsfs}

%Inserimento pseudocodice
\usepackage{algorithm}
\usepackage{algpseudocode}

%Citazione frasi
\usepackage{csquotes}

%Inserimento di Lorem ipsum nel testo
\usepackage{lipsum}

%\usepackage{lmodern}
\usepackage{pifont}


%Nuovi comandi
%*****************
%* nuovi comandi *
%*****************

\newcommand{\abs}[1]{\left|#1\right|}                               % modulo
\newcommand{\dato}{\left|\right.}                                   % probabilit\`{a} condizionata
\newcommand{\fun}[1]{\mathrm{#1}}                                   % stile funzione
\newcommand{\imp}{\;\;\Longrightarrow\;\;}                          % implicazione
\newcommand{\norma}[1]{\left\| #1 \right\|}                         % norma
\newcommand{\prob}[1]{\mathrm{P}\!\left[#1\right]}                  % probabilit\`{a}
\newcommand{\expect}[1]{\mathrm{E}\!\left[#1\right]}                % aspettazione
\newcommand{\sse}{\;\;\Longleftrightarrow\;\;}                      % se e solo se
\newcommand{\vect}[1]{{\boldsymbol{\mathrm{#1}}}}                   % stile vettore
\newcommand{\real}[1]{\fun{Re}\left[#1\right]}                      % parte reale
\newcommand{\imag}[1]{\fun{Im}\left[#1\right]}                      % parte immaginaria
\newcommand{\Dim}[1]{\fun{dim}\left[#1\right]}                      % dimensione di una matrice
\newcommand{\Det}[1]{\fun{det}\left[#1\right]}                      % determinante di una matrice
\newcommand{\Ker}[1]{\fun{ker}\left[#1\right]}                      % ker di una matrice
\newcommand{\rango}[1]{\fun{rango}\left[#1\right]}                  % rango di una matrice
\newcommand{\scalare}[2]{\left\langle #1, #2 \right\rangle}         % prodotto scalare
\newcommand{\blbrace}{\left  \lbrace}                               % parentesi graffa sinistra grande
\newcommand{\brbrace}{\right \rbrace}                               % parentesi graffa destra grande
\newcommand{\sinc}{\fun{sinc}}                                      % sinc
\newcommand{\rect}{\fun{rect}}                                      % rect
\newcommand{\rcos}{\fun{rcos}}                                      % rcos
\newcommand{\sgn}{\fun{sgn}}                                        % sgn
\newcommand{\N}{\mathbb{N}}                                         % insieme dei numeri naturali
\newcommand{\Z}{\mathbb{Z}}                                         % insieme dei numeri interi
\newcommand{\Q}{\mathbb{Q}}                                         % insieme dei numeri razionali
\newcommand{\R}{\mathbb{R}}                                         % insieme dei numeri reali
\newcommand{\C}{\mathbb{C}}                                         % insieme dei numeri complessi
\newcommand{\seq}[2][n]{#2_{0}, #2_{1}, \ldots, \, #2_{#1}}         % sequenza
\newcommand{\Span}[2][n]
{\fun{span} \blbrace #2_{1}, #2_{2}, \ldots, \, #2_{#1} \brbrace}   % spazio generato
\newcommand{\ddt}{\frac{\fun{d}}{\fun{dt}}}                         % derivata
\newcommand{\Div}[2]{#1 \; \mid \; #2}                              % divide
\newcommand{\MCD}[2]{\fun{MCD}\(#1, #2\)}                           % massimo comun divisore
\newcommand{\mcm}[2]{\fun{mcm}\(#1, #2\)}                           % minimo comune multiplo
\newcommand{\goodgap}{
                      \hspace{\subfigtopskip}
                      \hspace{\subfigbottomskip}
                     }                                              % interlinea opportuna per le sottofigure
\newcommand{\eng}[1]{\emph{#1}}                                   % inglese
\newcommand{\virg}[1]{``#1"}                                        % fa una citazione tra virgolette
%\newcommand{\unit}[2]($\frac{\text{#1}}{\text{#2}}$)                % unit\`{a} di misura
\newcommand{\textttvar}[1]{{\ttvar #1}}

\definecolor{gray}{gray}{0.9}
\newcommand{\listato}[1]{\lstset{language=#1, numbers=left, numberstyle=\tiny, stepnumber=2, numbersep=5pt, numberblanklines=false, xleftmargin=5pt, captionpos=b, stringstyle=\ttfamily, columns=flexible, showstringspaces=false, tabsize=2, frame=single, framerule=0pt, backgroundcolor=\color{gray}, basicstyle=\small}}

%****************************
%* ridefinizioni di comandi *
%****************************

\renewcommand{\(}{\left(}                                     % parentesi tonda sinistra grande
\renewcommand{\)}{\right)}                                    % parentesi tonda destra grande
\renewcommand{\[}{\left[}                                     % parentesi quadra sinistra grande
\renewcommand{\]}{\right]}                                    % parentesi quadra destra grande
\renewcommand{\exp}[1]{\fun{e}^{#1}}                          % esponenziale
\renewcommand{\gcd}[2]{\fun{gcd}\(#1, #2\)}                   % massimo comun divisore

\renewcommand{\lstlistingname}{Listing}
\renewcommand{\lstlistlistingname}{Elenco dei listati codice}

\makeatletter
\def\BState{\State\hskip-\ALG@thistlm}
\makeatother

%Definizione input e output pseudocodice
\renewcommand{\algorithmicrequire}{\textbf{Input:}}
\renewcommand{\algorithmicensure}{\textbf{Output:}}
 
%"Nome" della bibliografia
\addto{\captionsenglish}{%
  \renewcommand{\bibname}{References}
}

%Impostazioni dei margini, definizione di colori e di stili vari
\newfont{\ttvar}{cmvtt10 scaled 1200}     % nuovo carattere tipi courier a spaziatura variabile per le dimostrazioni

% margini senato
\textwidth       =  14.50 cm            % larghezza 21 cm - 4 cm (sinistro) - 2.5 (destro)
\textheight      =  23.10 cm            % altezza 29.7 cm - 3 cm (superiore) - 2 (inferiore)
\topmargin       =   0.00 cm            % margine superiore 3 cm diminuito di 1 inch
\oddsidemargin   =   1.46 cm            % margine sinistro 4 cm diminuito di 1 inch
\evensidemargin  =  -0.04 cm            % margine destro 2.5 cm diminuito di un inch

\setlength{\headsep}{1.0cm}
\setlength{\footskip}{1.0cm}
\parindent = 0.7cm
\captionmargin = 0.7cm

%Interlinea 1.5
\onehalfspacing

% stile pagina
\pagestyle{fancy}
\renewcommand{\chaptermark}[1]{\markboth{\chaptername\ \thechapter.\ #1 }{}}
\renewcommand{\sectionmark}[1]{\markright{\thesection\ #1}{}}
\fancyhead{}
\fancyhead[LE,RO]{\sffamily \thepage}
\fancyhead[RE]{\sffamily \leftmark}
\fancyhead[LO]{\sffamily \rightmark}
\fancyfoot{}

% ridefinisco lo stile plain
\fancypagestyle{plain}{ \fancyhead{} \fancyfoot{}
\fancyfoot[C]{\sffamily \thepage}
\renewcommand{\headrulewidth}{0pt}}

% stile per i titoli
\allsectionsfont{\sffamily \raggedright}

% definisco i colori
\definecolor{codegreen}{rgb}{0,0.6,0}
\definecolor{codegray}{rgb}{0.5,0.5,0.5}
\definecolor{codepurple}{rgb}{0.58,0,0.82}
\definecolor{backcolour}{rgb}{0.95,0.95,0.92}
\definecolor{backcolourwhite}{rgb}{1,1,1}

%definisco stile listati di codice
\lstdefinestyle{mystyle}{
    %backgroundcolor=\color{backcolour},   
    commentstyle=\color{codegreen},
    keywordstyle=\color{magenta},
    numberstyle=\tiny\color{codegray},
    stringstyle=\color{codepurple},
    basicstyle=\ttfamily\footnotesize,
    breakatwhitespace=false,         
    breaklines=true,                 
    captionpos=b,                    
    keepspaces=true,                 
    numbers=left,                    
    numbersep=5pt,                  
    showspaces=false,                
    showstringspaces=false,
    showtabs=false,                  
    tabsize=2,
    xleftmargin=.07\textwidth, %xrightmargin=.2\textwidth
}

\lstset{style=mystyle}
\lstset{emph={RandomForestClassifier, RandomizedSearchCV, GridSearchCV},emphstyle=\underbar}



\begin{document}
\begin{titlepage}
\begin{center}
%\includegraphics[scale=0.1]{images/logo.png}\\

%Per il frontespizio del dipartimenti di Ing. dell'Informazione commentare le riga precedente e decommentare la successiva
\includegraphics[scale=0.2]{images/logo_unipd.png} \hfill \includegraphics[scale=0.2]{images/logo_dei.png}\\
\vspace{0.8cm}
\textsc{\LARGE Universit\`{a} degli Studi di Padova}\\
\vspace{0.45cm}
\textsc{\large Dipartimento di Ingegneria dell'Informazione}\\
\vspace{0.4cm}
\textsc{\large Corso di Laurea Triennale in}\\
\textsc{\large Ingegneria Informatica}\\
\vfill
% Title
{ \LARGE \bfseries Analisi dei sensori coppia-forza per lo sviluppo di applicazioni industriali in ROS
}\\
\vfill
\textit{\large Relatore:} \hfill \textit{\large Laureando:}\\
\textsc{\large Prof. Stefano Ghidoni} \hfill \textsc{Andrea Stocco}\\
\textit{\large Correlatore:} \hfill \textsc{2009353}\\
\textsc{\large Matteo Terreran, PhD} \hfill \textit{}

\vfill
% Bottom of the page
{\large Anno Accademico 2022/2023}
\end{center}
\end{titlepage}

\thispagestyle{empty} %pagina bianca dopo il titolo
\cleardoublepage

\pagenumbering{Roman} %numerazione romana per l'indice, l'abstract e i ringraziamenti
\thispagestyle{empty}

\clearpage{\pagestyle{plain}\cleardoublepage}
%definisco il layout dell'abstract
\def\changemargin#1#2{\list{}{\rightmargin#2\leftmargin#1}\item[]}
\let\endchangemargin=\endlist

%Genero l'ambiente per l'abstract
\newcommand\summaryname{Abstract}
\newenvironment{Abstract}%
    {\begin{center}%
    \bfseries{\summaryname} \end{center}}

\begin{Abstract}
%\begin{changemargin}{1cm}{1cm}
    I sensori coppia-forza sono componenti fondamentali nei sistemi robotici in quanto forniscono dati sulle forze e i momenti 
    esterni applicati al robot. Questi dati possono essere utilizzati per applicazioni a supporto della collaborazione uomo-robot 
    e per l'automatizzazione di attivit\'{a} che richiedono elevata precisione. ROS (Robot Operating System) \'{e} un framework che 
    fornisce una vasta gamma di librerie e strumenti software per lo sviluppo di applicazioni robotiche. In questa tesi verranno 
    mostrate delle possibili applicazioni ROS in ambito industriale per i sensori coppia-forza, dopo averli validati attraverso 
    alcuni esperimenti.
%\end{changemargin}
\end{Abstract}

\clearpage{\pagestyle{plain}\cleardoublepage}
\tableofcontents %Indice

\clearpage{\pagestyle{plain}\cleardoublepage} %Numerazione araba per i capitoli
\pagenumbering{arabic}

\clearpage{\pagestyle{plain}\cleardoublepage} %Comando per iniziare il capitolo su pagina dispari
\chapter*{Introduzione} %Nome capitolo
\addcontentsline{toc}{chapter}{Introduzione} 
I robot manipolatori hanno rivoluzionato l'automazione industriale, permettendo lo svolgimento
di operazioni complesse in modo rapido, preciso e sicuro. 
Uno dei modelli pi\'{u} utilizzato \'{e} l'\textbf{UR5} di \textbf{Universal Robot}. Per via della sua flessibilit\'{a} 
ed efficienza si \'{e} scelto di utilizzarlo in questo studio. 
Un ruolo chiave nel controllo di questi robot viene assunto dai sensori coppia-forza,
che permettono di misurare e regolare la forza esercitata dal robot durante lo svolgimento delle proprie attivit\'{a}. 
La loro versatilit\'{a} li rende strumenti preziosi in molti campi tra cui: la \textbf{robotica industriale} e la \textbf{medicina}. 
Essi infatti consentono al robot di controllare la forza esercitata durante le operazioni di assemblaggio, levigatura, 
saldatura e manipolazione degli oggetti. Vengono inoltre utilizzati nella riabilitazione, per valutare la forza muscolare 
e i progressi del paziente.
% aggiungere foto levigatura e riabilitazione
Questi sensori sono in grado di convertire le forze e le coppie applicati ad essi in segnali elettrici che possono essere 
interpretati da altri dispositivi. 
Esistono varie tipologie di sensori coppia-forza, ognuna delle quali ha un diverso meccanismo di funzionamento. 
I sensori \textbf{piezoelettrici} sfruttano la propriet\'{a} di alcuni materiali (cristalli piezoelettrici) 
di generare una carica elettrica se sottoposti a deformazione meccanica. Tale variazione  
pu\'{o} essere misurata per determinare la forza o la coppia applicata.
Il sensore \textbf{FT 300-S} di \textbf{Robotiq} sfrutta questo principio di funzionamento ed \'{e} in grado di misurare 
forze e coppie lungo i sei gradi di libert\'{a} (x, y, z, roll, pitch, yaw). 
\'{E} stato necessario installarlo manualmente sull'UR5 perch\'{e}, a differenza di altri, non possiede sensori 
coppia-forza integrati. 
In questa tesi verr\'{a} presentata l'implementazione di un sistema di controllo della forza per l'UR5 utilizzando i dati
forniti dall'FT 300-S e il framework di sviluppo \textbf{ROS (Robot Operating System)}.
ROS \'{e} ampiamente utilizzato dalla comunit\'{a} informatica perch\'{e} fornisce strumenti e librerie 
per il controllo e la comunicazione tra le componenti di un sistema robotico.
Alcuni test per la valutazione delle prestazioni del sensore in termini di reattivit\'{a} e precisione verranno descritti nel Capitolo 
\ref{chapter:chapter4}.
Nel Capitolo \ref{chapter:chapter5}, invece, verranno presentate delle possibili applicazioni volte a dimostrare l'efficacia
di tali sensori per lo svolgimento di attivit\'{a} industriali, come il \textbf{pick and place} e il \textbf{trasporto collaborativo}.
 %File in cui verrà scritto il capitolo

\clearpage{\pagestyle{plain}\cleardoublepage} %Comando per iniziare il capitolo su pagina dispari
\chapter{ROS} %Nome capitolo
\label{chapter:chapter1} %Label per creare riferimenti al capitolo
ROS (Robot Operating System) \'{e} un framework open-source di Python e C++ per lo sviluppo di applicazioni robotiche.
Ne esistono varie versioni, quella raccomandata e utilizzata in questa tesi \'{e} ROS Noetic per Ubuntu Focal 20.04.
Questo capitolo ci fornir\'{a} una panoramica esaustiva del Robot Operating System e dell'ambiente di esecuzione associato 
\cite{ros}.

\section{Workspace}
Per poter eseguire codice ROS abbiamo bisogno di un ambiente che permetta l'organizzazione e l'utilizzo di tutti i nostri pacchetti.
Catkin \'{e} il sistema di compilazione ufficiale di ROS che ci consente di creare un workspace per organizzare 
e gestire le nostre applicazioni.
I termini `pacchetto' e `applicazione' sono interscambiabili e possono essere utilizzati in modo equivalente (una definizione pi\'{u} 
rigorosa verr\'{a} data nel corso di questo capitolo).
Una volta creato il workspace catkin \cite{catkin_ws}, saremo in grado di compilare il codice sorgente contenuto all'interno 
dei nostri pacchetti.


\section{Nodo}
Un \textbf{nodo} \'{e} un eseguibile che sfrutta ROS per comunicare con altri nodi.
Quando viene lanciato il comando \verb|catkin_make|, ogni file sorgente in ogni pacchetto viene compilato 
dando origine ad un nodo.
Impropriamente, si potrebbe quindi dire che un pacchetto \'{e} un insieme di nodi riguardanti la stessa applicazione.
Per poter eseguire un nodo \'{e} sufficiente eseguire il comando \\
\verb|rosrun <nome_pacchetto> <nome_nodo>|. 
Prima di eseguirne uno \'{e} necessario, tuttavia, avviare un ROS Master.
Lo scopo principale di un \textbf{ROS Master} \'{e} quello di consentire ai singoli nodi di localizzarsi a vicenda. Una volta 
fatto partire, essi potranno comunicare tra loro attraverso topic o service.
Per eseguire un ROS Master sar\'{a} sufficiente eseguire il comando \verb|roscore| in un altro terminale.


\section{Topic}
I topic sono dei canali di comunicazione unidirezionali che consentono lo scambio di informazioni tra nodi sottoforma di 
messaggi. Un nodo che pubblica messaggi su un topic viene chiamato publisher, mentre un nodo che legge i messaggi da un topic 
viene chiamato subscriber.
\begin{figure}[H]
    \centering
    \includegraphics*[width=0.75\textwidth]{images/topic_graph.png}
    \caption{Schema di comunicazione}
    \label{fig:topic_graph}
\end{figure}
Per pubblicare un messaggio su un topic bisogna utilizzare l'apposita funzione \verb|publish()| passandole come parametro 
il messaggio che vogliamo pubblicare.
A discapito del nome, questa funzione non pubblica effettivamente il messaggio, ma lo mette in una coda d'attesa (la cui 
dimensione viene specificata quando viene istanziato il publisher).
Un thread separato si occupa di inviare effettivamente il messaggio al topic per renderlo visibile a tutti i nodi subscriber 
connessi.
Se il numero di messaggi in coda supera la dimensione di essa, i messaggi pi\'{u} vecchi verranno cancellati per fare spazio 
a quelli pi\'{u} recenti. Quando arriva un nuovo messaggio al subscriber, esso viene salvato 
in una coda d'attesa (stesso funzionamento di quella del publisher) fino a quando ROS non d\'{a} la possibilit\'{a} al nodo 
di eseguire la funzione di callback. Tale funzione \'{e} definita dall'utente e si occupa di processare il messaggio ricevuto. 
ROS eseguir\'{a} una callback solo quando gli daremo il permesso di farlo. Ci sono due modi per fare ci\'{o}:
\begin{itemize}
    \item \verb|ros::spinOnce()| chiede a ROS di eseguire tutte le callback in sospeso e poi ci restituisce il controllo
    \item \verb|ros::spin()| chiede a ROS di attendere e di eseguire tutte le callback in sospeso fino a quando il nodo non 
          viene spento. \'{E} equivalente a:
          \begin{verbatim}
            while (ros::ok()) {
                ros::spinOnce();
            }
          \end{verbatim} 
          \verb|ros::ok()| ritorna 0 quando: 
          \begin{itemize}
            \item il nodo viene spento attraverso un SIGINT (Ctrl-C) oppure dalla chiamata di \verb|ros::shutdown()| in 
                  un altro punto del codice
            \item un altro nodo con lo stesso nome viene eseguito
          \end{itemize}
\end{itemize}
In altre parole \verb|ros::spin()| vincola il nodo a rimanere sempre e solo in attesa di leggere nuovi messaggi. 
Se il nodo non deve solamente eseguire le callback, allora un loop con \verb|ros::spinOnce()| \'{e} la scelta corretta.
(Aggiungere link ad un mio pacchetto di esempio???)
 %File in cui verrà scritto il capitolo

\clearpage{\pagestyle{plain}\cleardoublepage} %Comando per iniziare il capitolo su pagina dispari
\chapter{Ambiente di lavoro} %Nome capitolo
\label{chapter:chapter2} %Label per creare riferimenti al capitolo
Il sensore FT300-S di Robotiq \'{e} stato installato all'estremit\'{a} dell'UR5 e collegato alla control box per ricevere 
l'alimentazione necessaria. 
\begin{figure}[H]
    \centering
    \includegraphics*[width=0.5\textwidth]{images/ft.png}
    \caption{FT300-S}
    \label{fig:ft}
\end{figure}
\'{E} importante notare come la sua presenza non precluda la possibilit\'{a} di installazione di un end effector, 
che pu\'{o} essere facilmente posizionato `al di sopra' del sensore. 
L'FT300-S \'{e} in grado di rilevare forze e torsioni nel range di $\pm 300 N$ e $\pm 30 Nm$ rispettivamente. 
Le misurazioni del sensore hanno un rumore di fondo intrinseco, \'{e} quindi necessario scartare tutti i dati al di sotto delle 
soglie consigliate nel manuale \cite{ft_sensor} in quanto non attendibili. 
Per collegare il sensore al PC sono state provate due alternative: 
\begin{itemize}
    \item collegamento via USB tra sensore e control box e via ethernet tra control box e computer
    \item collegamento diretto via USB tra sensore e computer
\end{itemize}

\subsection{Collegamento via USB tra sensore e control box e via ethernet tra control box e computer}
\begin{figure}[H]
    \centering
    \includegraphics*[width=0.1\textwidth]{images/ft-cbox-pc.png}
    \caption{Schema collegamento}
    \label{fig:ft-cbox-pc}
\end{figure}
bla bla bla

\subsection{Collegamento diretto via USB tra sensore e computer}
\begin{figure}[H]
    \centering
    \includegraphics*[width=0.1\textwidth]{images/ft-pc.png}
    \caption{Schema collegamento}
    \label{fig:ft-pc}
\end{figure}
bla bla bla

 %File in cui verrà scritto il capitolo

\clearpage{\pagestyle{plain}\cleardoublepage} %Comando per iniziare il capitolo su pagina dispari
\chapter{Validazione del sensore} %Nome capitolo
\label{chapter:chapter3} %Label per creare riferimenti al capitolo
%Il sensore FT300-S di Robotiq \'{e} stato installato all'estremit\'{a} dell'UR5 e collegato alla control box per ricevere 
l'alimentazione necessaria. 
\begin{figure}[H]
    \centering
    \includegraphics*[width=0.5\textwidth]{images/ft.png}
    \caption{FT300-S}
    \label{fig:ft}
\end{figure}
\'{E} importante notare come la sua presenza non precluda la possibilit\'{a} di installazione di un end effector, 
che pu\'{o} essere facilmente posizionato `al di sopra' del sensore. 
L'FT300-S \'{e} in grado di rilevare forze e torsioni nel range di $\pm 300 N$ e $\pm 30 Nm$ rispettivamente. 
Le misurazioni del sensore hanno un rumore di fondo intrinseco, \'{e} quindi necessario scartare tutti i dati al di sotto delle 
soglie consigliate nel manuale \cite{ft_sensor} in quanto non attendibili. 
Per collegare il sensore al PC sono state provate due alternative: 
\begin{itemize}
    \item collegamento via USB tra sensore e control box e via ethernet tra control box e computer
    \item collegamento diretto via USB tra sensore e computer
\end{itemize}

\subsection{Collegamento via USB tra sensore e control box e via ethernet tra control box e computer}
\begin{figure}[H]
    \centering
    \includegraphics*[width=0.1\textwidth]{images/ft-cbox-pc.png}
    \caption{Schema collegamento}
    \label{fig:ft-cbox-pc}
\end{figure}
bla bla bla

\subsection{Collegamento diretto via USB tra sensore e computer}
\begin{figure}[H]
    \centering
    \includegraphics*[width=0.1\textwidth]{images/ft-pc.png}
    \caption{Schema collegamento}
    \label{fig:ft-pc}
\end{figure}
bla bla bla

 %File in cui verrà scritto il capitolo

\clearpage{\pagestyle{plain}\cleardoublepage} %Comando per iniziare il capitolo su pagina dispari
\chapter{Applicazioni industriali} %Nome capitolo
\label{chapter:chapter4} %Label per creare riferimenti al capitolo
%Il sensore FT300-S di Robotiq \'{e} stato installato all'estremit\'{a} dell'UR5 e collegato alla control box per ricevere 
l'alimentazione necessaria. 
\begin{figure}[H]
    \centering
    \includegraphics*[width=0.5\textwidth]{images/ft.png}
    \caption{FT300-S}
    \label{fig:ft}
\end{figure}
\'{E} importante notare come la sua presenza non precluda la possibilit\'{a} di installazione di un end effector, 
che pu\'{o} essere facilmente posizionato `al di sopra' del sensore. 
L'FT300-S \'{e} in grado di rilevare forze e torsioni nel range di $\pm 300 N$ e $\pm 30 Nm$ rispettivamente. 
Le misurazioni del sensore hanno un rumore di fondo intrinseco, \'{e} quindi necessario scartare tutti i dati al di sotto delle 
soglie consigliate nel manuale \cite{ft_sensor} in quanto non attendibili. 
Per collegare il sensore al PC sono state provate due alternative: 
\begin{itemize}
    \item collegamento via USB tra sensore e control box e via ethernet tra control box e computer
    \item collegamento diretto via USB tra sensore e computer
\end{itemize}

\subsection{Collegamento via USB tra sensore e control box e via ethernet tra control box e computer}
\begin{figure}[H]
    \centering
    \includegraphics*[width=0.1\textwidth]{images/ft-cbox-pc.png}
    \caption{Schema collegamento}
    \label{fig:ft-cbox-pc}
\end{figure}
bla bla bla

\subsection{Collegamento diretto via USB tra sensore e computer}
\begin{figure}[H]
    \centering
    \includegraphics*[width=0.1\textwidth]{images/ft-pc.png}
    \caption{Schema collegamento}
    \label{fig:ft-pc}
\end{figure}
bla bla bla

 %File in cui verrà scritto il capitolo

\clearpage{\pagestyle{plain}\cleardoublepage} %Comando per iniziare il capitolo su pagina dispari
\chapter{Conclusioni} %Nome capitolo
\label{chapter:chapter5} %Label per creare riferimenti al capitolo
%Il sensore FT300-S di Robotiq \'{e} stato installato all'estremit\'{a} dell'UR5 e collegato alla control box per ricevere 
l'alimentazione necessaria. 
\begin{figure}[H]
    \centering
    \includegraphics*[width=0.5\textwidth]{images/ft.png}
    \caption{FT300-S}
    \label{fig:ft}
\end{figure}
\'{E} importante notare come la sua presenza non precluda la possibilit\'{a} di installazione di un end effector, 
che pu\'{o} essere facilmente posizionato `al di sopra' del sensore. 
L'FT300-S \'{e} in grado di rilevare forze e torsioni nel range di $\pm 300 N$ e $\pm 30 Nm$ rispettivamente. 
Le misurazioni del sensore hanno un rumore di fondo intrinseco, \'{e} quindi necessario scartare tutti i dati al di sotto delle 
soglie consigliate nel manuale \cite{ft_sensor} in quanto non attendibili. 
Per collegare il sensore al PC sono state provate due alternative: 
\begin{itemize}
    \item collegamento via USB tra sensore e control box e via ethernet tra control box e computer
    \item collegamento diretto via USB tra sensore e computer
\end{itemize}

\subsection{Collegamento via USB tra sensore e control box e via ethernet tra control box e computer}
\begin{figure}[H]
    \centering
    \includegraphics*[width=0.1\textwidth]{images/ft-cbox-pc.png}
    \caption{Schema collegamento}
    \label{fig:ft-cbox-pc}
\end{figure}
bla bla bla

\subsection{Collegamento diretto via USB tra sensore e computer}
\begin{figure}[H]
    \centering
    \includegraphics*[width=0.1\textwidth]{images/ft-pc.png}
    \caption{Schema collegamento}
    \label{fig:ft-pc}
\end{figure}
bla bla bla

 %File in cui verrà scritto il capitolo

\clearpage{\pagestyle{plain}\cleardoublepage}
\begin{thebibliography}{99}
\addcontentsline{toc}{chapter}{Bibliografia} 

\bibitem{ros} \textit{Quigley, Morgan, et al. "ROS: an open-source Robot Operating System." ICRA workshop on open source software. 
Vol. 3. No. 3.2. 2009.}

\bibitem{ros_tutorial} \textit{ROS Tutorials}, \url{http://wiki.ros.org/ROS/Tutorials}.

\bibitem{catkin_ws} \textit{Catkin Workspace}, \url{http://wiki.ros.org/ROS/Tutorials/InstallingandConfiguringROSEnvironment}.

\bibitem{yt_tutorial} \textit{YouTube Tutorial}, \url{https://www.youtube.com/playlist?list=PLLSegLrePWgIbIrA4iehUQ-impvIXdd9Q}.

\bibitem{srv_topic_example} \textit{Esempio topic e service}, \url{https://github.com/andreastocco01/ros/tree/main}.

\bibitem{service} \textit{Creating ROS msg and srv} \url{http://wiki.ros.org/ROS/Tutorials/CreatingMsgAndSrv}.

\bibitem{ur5} \textit{Download material for UR5} \url{https://www.universal-robots.com/cb3/}.

\bibitem{ft_sensor} \textit{Download FT300-S manual} \url{https://robotiq.com/support/ft-300-force-torque-sensor}.

\bibitem{ft_driver} \textit{FT300-S driver} \url{https://github.com/andreastocco01/ft300_driver}.

\bibitem{robotiq_repo} \textit{Robotiq maintained repo} \url{https://github.com/TAMS-Group/robotiq}.

\bibitem{moveit_tutorial} \textit{MoveIt Tutorial} \url{https://ros-planning.github.io/moveit_tutorials/}.

\bibitem{environment_setup} \textit{Ambiente di lavoro} \url{https://github.com/andreastocco01/environment_setup}.
\end{thebibliography}
\end{document}